\section{The divisor obstruction}\label{sec:divisor}

\subsection{Statement and proof}

\begin{lemma}[Restriction from an ambient projective space]\label{lem:ambient}
Let $Y$ be a smooth projective variety, let $i\colon Y\to\PP^{N-1}$ be a
morphism defined by a linear system, let $h\in H^{2}(\PP^{N-1},\QQ)$ be the
hyperplane class, and let $\cF\in D^{b}(\PP^{N-1})$. Then
\[
  \ch\bigl(Li^{*}\cF\bigr)\in\QQ\bigl[\,i^{*}h\,\bigr]\subset H^{*}(Y,\QQ).
\]
\end{lemma}

\begin{proof}
The Chern character commutes with derived pullback:
$\ch(Li^{*}\cF)=i^{*}\ch(\cF)$. This is standard for a locally free sheaf,
where it follows from naturality of Chern classes, and it extends to
$D^{b}(\PP^{N-1})$ because every object there has a finite resolution by
direct sums of line bundles $\cO(m)$, by Beilinson's theorem, and the Chern
character is additive on distinguished triangles.%
\footnote{Beilinson's theorem gives more than is used here: every object of
$D^{b}(\PP^{N-1})$ is quasi-isomorphic to a finite complex whose terms are
direct sums of $\cO(-m)$ with $0\le m\le N-1$. The proof needs only that
$K_{0}(\PP^{N-1})$ is generated by the classes of line bundles, which already
follows from the Koszul resolutions of the structure sheaves of linear
subspaces; \Cref{rem:otherambient} records which other ambient varieties
behave in the same way.} The cohomology ring of
$\PP^{N-1}$ is $\QQ[h]/(h^{N})$, so $\ch(\cF)$ is a polynomial in $h$ with
rational coefficients, and applying $i^{*}$, a ring homomorphism, gives a
polynomial in $i^{*}h$.
\end{proof}

\begin{lemma}[The hyperplane class of a linear system]\label{lem:hyperclass}
Let $A$ be an abelian variety and $i\colon\Av\to\PP^{N-1}$ a morphism defined
by a linear system $|M|$ with $M\in\Pic(\Av)$. Then
$i^{*}h=c_{1}(M)\in\NS(\Av)\otimes\QQ$.
\end{lemma}

\begin{proof}
A morphism to $\PP^{N-1}$ defined by a linear system $|M|$, or by a subsystem
of it, satisfies $i^{*}\cO(1)\cong M$ by construction, whence
$i^{*}h=c_{1}(i^{*}\cO(1))=c_{1}(M)$.
\end{proof}

\begin{theorem}[\Cref{thm:divisor}]\label{thm:divisor-body}
Let $A$ be a principally polarised complex abelian variety of dimension $g$
with $\NS(\Av)\otimes\QQ=\QQ\cdot\theta'$, let $i\colon\Av\to\PP^{N-1}$ be any
morphism defined by a linear system, and let $\cF$ be any object of
$D^{b}(\PP^{N-1})$. Put $\cE=\Phi_{\cP}(Li^{*}\cF)$. Then
\[
  \ch_{k}(\cE)\in\QQ\cdot\theta^{k}\qquad\text{for every }k\ge0 .
\]
\end{theorem}

\begin{proof}
By \Cref{lem:ambient}, $\ch(Li^{*}\cF)\in\QQ[i^{*}h]$, and by
\Cref{lem:hyperclass}, $i^{*}h=c_{1}(M)$ for the line bundle $M$ defining $i$.
By hypothesis $\NS(\Av)\otimes\QQ=\QQ\theta'$, so $c_{1}(M)=m\theta'$ for some
$m\in\QQ$, and therefore $\ch(Li^{*}\cF)\in\QQ[\theta']$. By
\Cref{cor:stability}, $\Phi_{\cP}$ carries $\QQ[\theta']$ into $\QQ[\theta]$,
so $\ch(\cE)\in\QQ[\theta]$ by \Cref{prop:grr}. Taking the component in
$H^{2k}$ and using that $\QQ[\theta]\cap H^{2k}(A,\QQ)=\QQ\theta^{k}$ gives the
statement.
\end{proof}

\begin{remark}[On the hypothesis on $\NS$]\label{rem:NShyp}
The hypothesis $\NS(\Av)\otimes\QQ=\QQ\theta'$ holds at the very general point
of any positive-dimensional family of abelian varieties of Weil type with
$\End^{0}(A)=\iota(K)$, which is the situation in which the Hodge conjecture
for Weil classes is open; see \Cref{lem:genericNS}. Without it the conclusion
must be stated as $\ch_{k}(\cE)\in D^{2k}(\Av)$ transported to $A$, that is,
$\ch_{k}(\cE)$ lies in the image under $\Phi_{\cP}$ of the divisor subring of
$\Av$. \Cref{thm:character} covers the general case directly and without any
hypothesis on $\NS$, which is one reason to have two proofs.
\end{remark}

\begin{lemma}\label{lem:genericNS}
Let $\mathcal{A}\to S$ be the universal family of abelian varieties of
$(K,-1,n)$-Weil type over the corresponding moduli space, with $n\ge2$. At a
very general point $s\in S$ one has $\End^{0}(A_{s})=\iota(K)$ and
$\NS(A_{s})\otimes\QQ=\QQ\eta_{s}$, of dimension one.
\end{lemma}

\begin{proof}
The locus in $S$ where $\End^{0}$ is strictly larger than $\iota(K)$ is a
countable union of proper closed subvarieties, since each extra endomorphism
imposes a nontrivial algebraic condition on the period point and there are
countably many possible endomorphism algebras; its complement is the very
general locus. At such a point the Hodge group is the full special unitary
group $\mathrm{SU}(V,\psi)$ of the $K$-hermitian form $\psi$ attached to the
polarisation, by Weil's computation of the Hodge group of a general abelian
variety of Weil type. The space $\Hdg^{1}(A_{s})=\NS(A_{s})\otimes\QQ$ is the
space of Hodge-group invariants in $H^{2}(A_{s},\QQ)$. By \Cref{prop:grading}
in degree two,
$H^{2}\otimes\CC=\bigwedge^{2}V_{+}\oplus(V_{+}\otimes V_{-})\oplus\bigwedge^{2}V_{-}$,
and $\mathrm{SU}(V,\psi)$ acts on $V_{+}\otimes V_{-}$ with a one-dimensional
space of invariants, spanned by $\psi$ itself, and on
$\bigwedge^{2}V_{\pm}$ with no invariants for $2n\ge3$, since
$\bigwedge^{2}$ of the standard representation of $\mathrm{SL}_{2n}$ has no
invariant vector unless $2n=2$. Hence the invariants are one-dimensional and
spanned by $\eta_{s}$.
\end{proof}

\subsection{Consequences for Weil classes}

\begin{proof}[Proof of \Cref{cor:neverweil}]
Let $g=2n$ and let $(A,\iota,\eta)$ be of $(K,-1,n)$-Weil type with
$\NS(A)\otimes\QQ=\QQ\eta$. By \Cref{thm:divisor-body},
$\ch_{n}(\cE)\in\QQ\cdot\theta^{n}=\QQ\cdot\eta^{n}\subset D^{2n}(A)$. By
\Cref{prop:notpoly}, $\HW(A)\cap D^{2n}(A)=0$. So if $\ch_{n}(\cE)$ is a Weil
class then it lies in $\HW(A)\cap D^{2n}(A)=0$.
\end{proof}

\begin{proof}[Second proof of \Cref{cor:neverweil}, via characters]
By \Cref{lem:NSchar} the class $\eta^{n}$ is a $\Kx$-eigenclass for the
character $\Nm^{n}$, hence so is every element of $\QQ\eta^{n}$. By
\Cref{thm:divisor-body}, $\ch_{n}(\cE)\in\QQ\eta^{n}$, so $\ch_{n}(\cE)$ is an
eigenclass for $\Nm^{n}$. By \Cref{cor:equivariant} it is a Weil class only if
it vanishes.
\end{proof}

The second proof is longer as written but it is the one that generalises,
because it does not use the hypothesis on $\NS$, only the equivariance. We
record the version that does not mention the ambient space at all.

\begin{theorem}[Equivariance suffices]\label{thm:equivsuffices}
Let $(A,\iota,\eta)$ be of $(K,-1,n)$-Weil type and let $\cE\in D^{b}(A)$ be
an object such that for every $\tau$ in the order
$\iota^{-1}(\End(A))\subset K$ one has
\[
  \iota(\tau)^{*}\ch_{n}(\cE)=\Nm_{K/\QQ}(\tau)^{n}\,\ch_{n}(\cE).
\]
Then $\ch_{n}(\cE)$ is a Weil class only if $\ch_{n}(\cE)=0$.
\end{theorem}

\begin{proof}
The hypothesis says $\ch_{n}(\cE)$ lies in the isotypic component of $\Nm^{n}$
for the action of the order, and this component is the same as the one for the
full group $\Kx$ because the order spans $K$ over $\QQ$ and the action is
$\QQ$-linear. Apply \Cref{cor:equivariant}.
\end{proof}

\begin{remark}[What the two hypotheses cost]\label{rem:twohyp}
\Cref{thm:divisor-body} asks that the object be restricted from an ambient
projective space and imposes a condition on $\NS(\Av)$.
\Cref{thm:equivsuffices} asks only that the resulting class be equivariant
with the norm character. The second hypothesis is weaker and is satisfied
automatically by any construction whose input is $K$-invariant, which is the
case for every secant construction we are aware of, since secant varieties of
a $K$-invariant subvariety are $K$-invariant. The reason to state both is
that the first is checkable by inspection of the construction and the second
requires an equivariance computation.
\end{remark}

\begin{figure}[htbp]
\centering
  \includegraphics[width=0.96\textwidth]{fig_chernline}
\caption{The linear algebra behind \Cref{thm:divisor}, drawn inside
$H^{2n}(A,\QQ)$. The Weil line $W(A)$ is a plane, free of rank one over $K$
and of rank two over $\QQ$, and $\Kx$ acts on it through $\tau^{2n}$, so the
orbit of any nonzero class is a spiral that spans the plane and lies on no
line in it. The transverse axis carries $\eta^{n}$, on which $\Kx$ acts
through the norm. A construction returning a single eigenclass returns a
line, and by \Cref{prop:orbitplane} no line inside $W(A)$ is stable.}
\label{fig:route}
\end{figure}
